% --- Archon physics formalization source begin ---
% archon:physics
% archon:covers PhyXMiniProblems/problem_phyx_mini_0087.lean
% archon:source-report reports/phyx_mini/problem_phyx_mini_0087.source.json

\chapter{Physics problem phyx\_mini\_0087}
\label{ch:PhyXMiniProblems_problem_phyx_mini_0087}

\paragraph{Problem source.}
\#\# Physical scenario

In figure, two isotropic point sources \textbackslash{}( S\_1 \textbackslash{}) and \textbackslash{}( S\_2 \textbackslash{}) emit light at wavelength \textbackslash{}( \textbackslash{}lambda = 400\textbackslash{},\textbackslash{}text\{nm\} \textbackslash{}). Source \textbackslash{}( S\_1 \textbackslash{}) is located at \textbackslash{}( y = 640\textbackslash{},\textbackslash{}text\{nm\} \textbackslash{}); source \textbackslash{}( S\_2 \textbackslash{}) is located at \textbackslash{}( y = -640\textbackslash{},\textbackslash{}text\{nm\} \textbackslash{}). At point \textbackslash{}( P\_1 \textbackslash{}) (at \textbackslash{}( x = 720\textbackslash{},\textbackslash{}text\{nm\} \textbackslash{})), the wave from \textbackslash{}( S\_2 \textbackslash{}) arrives ahead of the wave from \textbackslash{}( S\_1 \textbackslash{}) by a phase difference of \textbackslash{}( 0.600\textbackslash{}pi\textbackslash{},\textbackslash{}text\{rad\} \textbackslash{}).

\#\# Question

What multiple of \textbackslash{}( \textbackslash{}lambda \textbackslash{}) gives the phase difference between the waves from the two sources as the waves arrive at point \textbackslash{}( P\_2 \textbackslash{}), which is located at \textbackslash{}( y = 720\textbackslash{},\textbackslash{}text\{nm\} \textbackslash{}).

\#\# Answer choices

- A: A: 2.60λ

- B: B: 2.70λ

- C: C: 2.90λ

- D: D: 3.20λ

\#\# Figure caption (auxiliary; use the image as primary evidence)

The image shows a two-dimensional coordinate system with labeled points. The horizontal axis is labeled "x" and the vertical axis is labeled "y." 

- There are four labeled points, two in black and two in red.
- The black points are labeled \textbackslash{}(P\_1\textbackslash{}) and \textbackslash{}(P\_2\textbackslash{}). \textbackslash{}(P\_1\textbackslash{}) is on the x-axis to the right, and \textbackslash{}(P\_2\textbackslash{}) is on the y-axis above the origin.
- The red points are labeled \textbackslash{}(S\_1\textbackslash{}) and \textbackslash{}(S\_2\textbackslash{}). \textbackslash{}(S\_1\textbackslash{}) and \textbackslash{}(S\_2\textbackslash{}) are on the y-axis, among which \textbackslash{}(S\_1\textbackslash{}) is above the x-axis and \textbackslash{}(S\_2\textbackslash{}) is below the x-axis.

The axes and points imply a relationship or coordinates system setup, indicating positions in a mathematical or geometrical context.

\#\# Dataset metadata

Wave Optics; Spatial Relation Reasoning, Physical Model Grounding Reasoning

\paragraph{Recorded answer/context.}
C: C: 2.90λ

\paragraph{Figure/image path.}
/root/proposal\_for\_physic/science-mango/phyx\_mini\_run/phyx\_data/test\_image/87.png

\paragraph{Formalization target.}
create a compiling Lean file with sorry bodies at `PhyXMiniProblems/problem\_phyx\_mini\_0087.lean`. The Lean declarations must preserve the physical quantities, dimensions or dimensional roles, figure labels, governing-law hypotheses, and final relation expressed by this problem.
Use Mathlib/Physlib names found through LeanExplore where available. If a physics API is missing, introduce faithful local abstractions rather than scalar placeholder aliases.

\begin{theorem}[Physics formalization target]
\label{thm:physics:phyx_mini_0087:target}
\lean{PhyXMiniProblems.ProblemPhyXMini0087.phaseDifferenceAtP2_eq_recordedAnswerC}
\uses{def:physics:phyx-mini-0087:phyxminiproblems-problemphyxmini0087-twopointsourcesetup, def:physics:phyx-mini-0087:phyxminiproblems-problemphyxmini0087-phaseleadinwavelengths, def:physics:phyx-mini-0087:phyxminiproblems-problemphyxmini0087-hasstatedfigurereadouts, def:physics:phyx-mini-0087:phyxminiproblems-problemphyxmini0087-hasstatedphasecalibrationatp1, def:physics:phyx-mini-0087:phyxminiproblems-problemphyxmini0087-satisfiesisotropicphasepropagation, def:physics:phyx-mini-0087:phyxminiproblems-problemphyxmini0087-answerphasedifferenceinwavelengths}
The assigned autoformalize agent should translate this physics problem into Lean declarations in the covered file, with theorem and lemma proof bodies written as `by sorry`.
\end{theorem}
\begin{proof}
This is an autoformalization task, not a proof task. Produce statements that can later be proved by the physics prover without weakening the physical model.
\end{proof}

% --- Archon declaration topology begin ---
\section{Declaration topology}
\begin{definition}[Dim Length]
  \label{def:physics:phyx-mini-0087:phyxminiproblems-problemphyxmini0087-dimlength}
  \lean{PhyXMiniProblems.ProblemPhyXMini0087.DimLength}
  A unit-aware physical length
\end{definition}
\begin{proof}
  This is the declared model definition of Dim Length.
\end{proof}
\begin{definition}[length Value In]
  \label{def:physics:phyx-mini-0087:phyxminiproblems-problemphyxmini0087-lengthvaluein}
  \lean{PhyXMiniProblems.ProblemPhyXMini0087.lengthValueIn}
  \uses{def:physics:phyx-mini-0087:phyxminiproblems-problemphyxmini0087-dimlength}
  Read a physical length as a real scalar in the specified length unit
\end{definition}
\begin{proof}
  This is the declared model definition of length Value In.
\end{proof}
\begin{definition}[nanometers Value]
  \label{def:physics:phyx-mini-0087:phyxminiproblems-problemphyxmini0087-nanometersvalue}
  \lean{PhyXMiniProblems.ProblemPhyXMini0087.nanometersValue}
  \uses{def:physics:phyx-mini-0087:phyxminiproblems-problemphyxmini0087-dimlength, def:physics:phyx-mini-0087:phyxminiproblems-problemphyxmini0087-lengthvaluein}
  The nanometre readout used for all coordinates and path lengths
\end{definition}
\begin{proof}
  This is the declared model definition of nanometers Value.
\end{proof}
\begin{definition}[Source Label]
  \label{def:physics:phyx-mini-0087:phyxminiproblems-problemphyxmini0087-sourcelabel}
  \lean{PhyXMiniProblems.ProblemPhyXMini0087.SourceLabel}
  Labels of the two isotropic point sources in the figure
\end{definition}
\begin{proof}
  This is the declared model definition of Source Label.
\end{proof}
\begin{definition}[Observation Label]
  \label{def:physics:phyx-mini-0087:phyxminiproblems-problemphyxmini0087-observationlabel}
  \lean{PhyXMiniProblems.ProblemPhyXMini0087.ObservationLabel}
  Labels of the two observation points in the figure
\end{definition}
\begin{proof}
  This is the declared model definition of Observation Label.
\end{proof}
\begin{definition}[Isotropic Point Source]
  \label{def:physics:phyx-mini-0087:phyxminiproblems-problemphyxmini0087-isotropicpointsource}
  \lean{PhyXMiniProblems.ProblemPhyXMini0087.IsotropicPointSource}
  One monochromatic isotropic point source. Its emission phase is an unwrapped radian phase at a common reference time; isotropic propagation is imposed separately by SatisfiesIsotropicPhasePropagation
\end{definition}
\begin{proof}
  This is the declared model definition of Isotropic Point Source.
\end{proof}
\begin{definition}[Two Point Source Setup]
  \label{def:physics:phyx-mini-0087:phyxminiproblems-problemphyxmini0087-twopointsourcesetup}
  \lean{PhyXMiniProblems.ProblemPhyXMini0087.TwoPointSourceSetup}
  \uses{def:physics:phyx-mini-0087:phyxminiproblems-problemphyxmini0087-dimlength, def:physics:phyx-mini-0087:phyxminiproblems-problemphyxmini0087-sourcelabel, def:physics:phyx-mini-0087:phyxminiproblems-problemphyxmini0087-observationlabel, def:physics:phyx-mini-0087:phyxminiproblems-problemphyxmini0087-isotropicpointsource}
  The two-source optical setup, including the unwrapped phase of each wave at each named observation point at one common observation time
\end{definition}
\begin{proof}
  This is the declared model definition of Two Point Source Setup.
\end{proof}
\begin{definition}[x Coordinate]
  \label{def:physics:phyx-mini-0087:phyxminiproblems-problemphyxmini0087-xcoordinate}
  \lean{PhyXMiniProblems.ProblemPhyXMini0087.xCoordinate}
  The x coordinate readout of a point in the chosen chart
\end{definition}
\begin{proof}
  This is the declared model definition of x Coordinate.
\end{proof}
\begin{definition}[y Coordinate]
  \label{def:physics:phyx-mini-0087:phyxminiproblems-problemphyxmini0087-ycoordinate}
  \lean{PhyXMiniProblems.ProblemPhyXMini0087.yCoordinate}
  The y coordinate readout of a point in the chosen chart
\end{definition}
\begin{proof}
  This is the declared model definition of y Coordinate.
\end{proof}
\begin{definition}[phase Lead Radians]
  \label{def:physics:phyx-mini-0087:phyxminiproblems-problemphyxmini0087-phaseleadradians}
  \lean{PhyXMiniProblems.ProblemPhyXMini0087.phaseLeadRadians}
  \uses{def:physics:phyx-mini-0087:phyxminiproblems-problemphyxmini0087-sourcelabel, def:physics:phyx-mini-0087:phyxminiproblems-problemphyxmini0087-observationlabel, def:physics:phyx-mini-0087:phyxminiproblems-problemphyxmini0087-twopointsourcesetup}
  Signed phase by which leading is ahead of trailing at an observation point, measured in unwrapped radians
\end{definition}
\begin{proof}
  This is the declared model definition of phase Lead Radians.
\end{proof}
\begin{definition}[phase Lead In Wavelengths]
  \label{def:physics:phyx-mini-0087:phyxminiproblems-problemphyxmini0087-phaseleadinwavelengths}
  \lean{PhyXMiniProblems.ProblemPhyXMini0087.phaseLeadInWavelengths}
  \uses{def:physics:phyx-mini-0087:phyxminiproblems-problemphyxmini0087-sourcelabel, def:physics:phyx-mini-0087:phyxminiproblems-problemphyxmini0087-observationlabel, def:physics:phyx-mini-0087:phyxminiproblems-problemphyxmini0087-twopointsourcesetup, def:physics:phyx-mini-0087:phyxminiproblems-problemphyxmini0087-phaseleadradians}
  The signed phase lead expressed as a number of wavelength cycles. This is a dimensionless multiple of λ: one wavelength of phase is 2π radians
\end{definition}
\begin{proof}
  This is the declared model definition of phase Lead In Wavelengths.
\end{proof}
\begin{definition}[Has Stated Figure Readouts]
  \label{def:physics:phyx-mini-0087:phyxminiproblems-problemphyxmini0087-hasstatedfigurereadouts}
  \lean{PhyXMiniProblems.ProblemPhyXMini0087.HasStatedFigureReadouts}
  \uses{def:physics:phyx-mini-0087:phyxminiproblems-problemphyxmini0087-nanometersvalue, def:physics:phyx-mini-0087:phyxminiproblems-problemphyxmini0087-twopointsourcesetup, def:physics:phyx-mini-0087:phyxminiproblems-problemphyxmini0087-xcoordinate, def:physics:phyx-mini-0087:phyxminiproblems-problemphyxmini0087-ycoordinate}
  The problem and primary-figure readouts, in the nanometre coordinate chart: λ = 400 nm, S₁ = (0, 640) nm and S₂ = (0, -640) nm, P₁ = (720, 0) nm and P₂ = (0, 720) nm. This predicate contains no arrival-phase conclusion at P₂
\end{definition}
\begin{proof}
  This is the declared model definition of Has Stated Figure Readouts.
\end{proof}
\begin{definition}[Has Stated Phase Calibration At P1]
  \label{def:physics:phyx-mini-0087:phyxminiproblems-problemphyxmini0087-hasstatedphasecalibrationatp1}
  \lean{PhyXMiniProblems.ProblemPhyXMini0087.HasStatedPhaseCalibrationAtP1}
  \uses{def:physics:phyx-mini-0087:phyxminiproblems-problemphyxmini0087-twopointsourcesetup, def:physics:phyx-mini-0087:phyxminiproblems-problemphyxmini0087-phaseleadradians}
  The measured calibration at P₁: the wave from S₂ is ahead of the wave from S₁ by 0.600π = 3π/5 radians. This is source data, not the requested phase relation at P₂
\end{definition}
\begin{proof}
  This is the declared model definition of Has Stated Phase Calibration At P1.
\end{proof}
\begin{definition}[Satisfies Isotropic Phase Propagation]
  \label{def:physics:phyx-mini-0087:phyxminiproblems-problemphyxmini0087-satisfiesisotropicphasepropagation}
  \lean{PhyXMiniProblems.ProblemPhyXMini0087.SatisfiesIsotropicPhasePropagation}
  \uses{def:physics:phyx-mini-0087:phyxminiproblems-problemphyxmini0087-nanometersvalue, def:physics:phyx-mini-0087:phyxminiproblems-problemphyxmini0087-twopointsourcesetup}
  Radial phase propagation for monochromatic isotropic point sources. At a fixed common time, propagation over path length r subtracts 2π r / λ from the source phase. Physlib's Space metric and the wavelength readout both use the figure's nanometre unit here. This law is uniform over both sources and both observation points; it does not assert the requested numerical phase difference at P₂
\end{definition}
\begin{proof}
  This is the declared model definition of Satisfies Isotropic Phase Propagation.
\end{proof}
\begin{definition}[Answer Choice]
  \label{def:physics:phyx-mini-0087:phyxminiproblems-problemphyxmini0087-answerchoice}
  \lean{PhyXMiniProblems.ProblemPhyXMini0087.AnswerChoice}
  Labels of the four answer choices printed in the problem
\end{definition}
\begin{proof}
  This is the declared model definition of Answer Choice.
\end{proof}
\begin{definition}[answer Phase Difference In Wavelengths]
  \label{def:physics:phyx-mini-0087:phyxminiproblems-problemphyxmini0087-answerphasedifferenceinwavelengths}
  \lean{PhyXMiniProblems.ProblemPhyXMini0087.answerPhaseDifferenceInWavelengths}
  \uses{def:physics:phyx-mini-0087:phyxminiproblems-problemphyxmini0087-answerchoice}
  Each answer as an exact dimensionless multiple of the wavelength λ
\end{definition}
\begin{proof}
  This is the declared model definition of answer Phase Difference In Wavelengths.
\end{proof}
\begin{lemma}[source paths to P1 are equal]
  \label{lem:physics:phyx-mini-0087:phyxminiproblems-problemphyxmini0087-source-paths-to-p1-are-equal}
  \lean{PhyXMiniProblems.ProblemPhyXMini0087.source_paths_to_P1_are_equal}
  \uses{def:physics:phyx-mini-0087:phyxminiproblems-problemphyxmini0087-twopointsourcesetup, def:physics:phyx-mini-0087:phyxminiproblems-problemphyxmini0087-hasstatedfigurereadouts}
  The primary-figure coordinates make the two source-to-P₁ path lengths equal. Thus the stated arrival-phase lead at P₁ calibrates the relative emission phase without a propagation-path correction
\end{lemma}
\begin{proof}
  The conclusion follows from the governing relations and figure readouts named in the statement of source paths to P1 are equal.
\end{proof}
\begin{lemma}[path difference to P2 in wavelengths]
  \label{lem:physics:phyx-mini-0087:phyxminiproblems-problemphyxmini0087-path-difference-to-p2-in-wavelengths}
  \lean{PhyXMiniProblems.ProblemPhyXMini0087.path_difference_to_P2_in_wavelengths}
  \uses{def:physics:phyx-mini-0087:phyxminiproblems-problemphyxmini0087-nanometersvalue, def:physics:phyx-mini-0087:phyxminiproblems-problemphyxmini0087-twopointsourcesetup, def:physics:phyx-mini-0087:phyxminiproblems-problemphyxmini0087-hasstatedfigurereadouts}
  At P₂, the path from S₂ is longer than the path from S₁ by 3.20λ. This is a geometric intermediate conclusion, not an assumed answer to the phase-difference question
\end{lemma}
\begin{proof}
  The conclusion follows from the governing relations and figure readouts named in the statement of path difference to P2 in wavelengths.
\end{proof}
% --- Archon declaration topology end ---
% --- Archon physics formalization source end ---
